\section{Justification}

Artificial Neural Networks (ANNs) represent one of the most relevant tools within the field of artificial intelligence due to their ability to learn from data and solve complex problems of classification, prediction, and pattern recognition.  
However, traditional machine learning methods present a critical limitation: when new data blocks are incorporated, the model tends to deteriorate in performance and forget part of the information previously acquired—a phenomenon known as catastrophic forgetting \cite{bullinaria2009}.

Despite recent advances in deep architectures and regularisation techniques, neural networks still lack enduring memory mechanisms that allow them to retain information over the long term without degrading their performance.  
In contrast, human beings have the ability to learn new tasks without significantly erasing previous knowledge, integrating new information progressively.  
Although this property is associated with the biological structure of the brain—particularly with processes such as synaptic consolidation—from a computational standpoint it is feasible to explore strategies that emulate this behaviour, optimising the storage and updating of knowledge in artificial models.

From a technical perspective, there are no theoretical limitations that prevent experimentation with new neural architectures, loss functions, or training schemes that promote progressive learning.  
The challenge lies in enabling the network to assimilate new data without overwriting previously acquired knowledge, thereby achieving a balance between the stability and plasticity of the system.  
Overcoming this challenge would enable the development of more robust and adaptive models capable of operating continuously in dynamic environments such as intelligent monitoring systems, robotics, or real-time applications where data are constantly updated.

In environments where incremental learning is not employed, each data incorporation requires retraining the entire model with all the available information.  
This process is costly in terms of computational time, energy consumption, and hardware requirements.  
Conversely, an incremental approach allows the model to process only the new data, updating its knowledge without repeating the entire training process.  
In addition to optimising resources, this characteristic can provide both economic and technological advantages over other methods, by reducing computational cost, development time, and algorithmic complexity in continuous learning applications.

There are currently various tools that facilitate the construction and training of neural networks.  
For the experimental development of this research, the Google Colab environment will be primarily employed, as it offers free infrastructure based on Graphics Processing Units (GPU), full integration with the Python programming language, and an interactive environment that supports fast and reproducible experimentation.

Likewise, the model implementation will be carried out using the PyTorch library, which is widely recognised for its flexibility and dynamism in handling computational graphs.  
PyTorch allows defining, modifying, and debugging network architectures intuitively during execution, thanks to its design based on dynamic tensors and its syntax closely aligned with pure Python.  
These characteristics facilitate the exploration of different training schemes, the customisation of loss functions, and the integration with scientific libraries such as NumPy and Matplotlib.

In this sense, PyTorch is considered a suitable tool for the development and evaluation of incremental learning models, as it provides a flexible environment that supports both experimentation and performance analysis under progressive data incorporation scenarios.  
Its use will allow the validation of the effects of incremental learning in mitigating catastrophic forgetting, contributing to the design of more efficient, sustainable, and human-like neural network architectures.
